%%%%%%%%%%%%%%%%%%%%%%%%%%%%%%%%%%%%%%%%%
%
% Generic Thesis Template
%
% This file is intentionally institution-agnostic and can be reused for
% different universities or departments by updating metadata and section content.
%
% The class and style files define formatting behavior. Keep them stable unless
% you intentionally need to customize LaTeX internals.
%
%%%%%%%%%%%%%%%%%%%%%%%%%%%%%%%%%%%%%%%%%


%----------------------------------------------------------------------------------------
%	PACKAGES AND DOCUMENT CONFIGURATIONS
%----------------------------------------------------------------------------------------
% class options (set in \documentclass[...]; order of options is irrelevant):
%   degree:       [wipro, bachelor, master, cas, sas]
%   language:     [english, german]
%   groupwork:    [groupwork] (extends declaration for two students)
%   ip:           [noip] (only for cas/sas if no intellectual property consent is required)
% 
% Examples for configuring the documentclass:
% \documentclass[german,master]{thesis}
% \documentclass[english,bachelor]{thesis}
% \documentclass[german,wipro,groupwork]{thesis}
% \documentclass[english,cas,noip]{thesis}
\documentclass[english,master]{thesis}


\usepackage{comment}                                  % having comment sections \begin{comment} \end{comment}
\usepackage{amsmath}							      % math package
\usepackage{amsfonts}							      % font package for math symbols
\usepackage{amssymb}							      % symbols package - definition of math symbols
\usepackage{listings}							      % package for code representation
\usepackage{graphicx}							      % for inclusion of image
\usepackage{subfig}								      % to arrange figures next to each other
\usepackage{float}								      % text style surrounding images
\usepackage[acronym]{glossaries}         		      % package for glossary
\usepackage{tikz}								      % used to place logos on title page
% \usepackage{gensymb}							      % for special characters such as °
\usepackage[a-1a]{pdfx}                               % Forces PDF/A-1a compliance for long-term archiving


\usepackage{multirow}
\usepackage{siunitx}
\usepackage{tabularx}
\usepackage{tikzscale}


\hypersetup{hidelinks}                                % hide red border in hyperlinks
\setcounter{tocdepth}{1}                              % hide subsections from TOC
\makenoidxglossaries      
\newacronym{org}{ORG}{Organization Name}
\newacronym{nn}{NN}{Neural Network}
                                      % include acronyms.txt file
\input{glossary}                                      % include glossary.txt file
\graphicspath{{figs/}}						          % set path of graphics folder

%----------------------------------------------------------------------------------------
%	INSTITUTION CUSTOMIZATION (GENERIC PLACEHOLDERS)
%----------------------------------------------------------------------------------------
\setUniName{University Name}
\setUniCountry{Country}
\setDepName{Department Name}
\setAcademicTitle{Degree Program Name}
\setPresentedTo{presented to Department Name}
\setPresentedToAcGrade{presented in consideration for the award of a degree}
\setBachelorThesis{Thesis}
\setMasterThesis{Thesis}
\setCasThesis{Thesis}
\setWiproThesis{Thesis}
\setConsSelf{I hereby declare that I have completed this work independently and that all sources, literature, and tools used are properly acknowledged and cited in accordance with academic integrity and applicable regulations.}
\setConsIp{Intellectual property and publication rights are subject to your institution's regulations.}


%----------------------------------------------------------------------------------------
%	PDF/A DOCUMENT COMPLIANCE
%----------------------------------------------------------------------------------------
\pdfcatalog{
  /StructTreeRoot <<                                  % Define the structure tree root for document tagging
    /Type /StructTreeRoot                             % Specify that this is a structure tree root
    /K []                                             % Placeholder for structure elements (empty for now)
  >>                
  /MarkInfo << /Marked true >>                        % Ensure the document is marked as tagged for accessibility
}


\begin{document}
%----------------------------------------------------------------------------------------
%	DOCUMENT INFORMATION
%----------------------------------------------------------------------------------------
\author{Your Full Name}                               % author name
\secondauthor{Second Author Name}                     % name of second author: only relevant for groupwork (groupwork flag must be set)
\city{Your City (Country)}                            % author's place of origin; 2nd place only required in the case of double authorship
\title{Your Thesis Title}                             % thesis title
\subtitle{\large A concise and descriptive subtitle}  % thesis subtitle

\date{2026}                                           % the year when the thesis was written (used in titlepage)
\defensedate{Month Day, 2026}                         % the date of the private defense
\defencelocation{Location}                            % location of defence
\extexpert{External Expert Name}                      % name of external expert
\indpartner{Industry Partner Name}                    % name of industry partner
\studyprogram{MSc ASE / MSc CSF}                      % name of study program

% jury, supervisor and dean are only relevant if acceptance sheet is enabled with the next line
% \addAcceptsheet
\jury{                                                % members of the jury
    \begin{itemize}
        \item Prof. Dr. Jury Member Name (President of the Jury);
        \item Prof. Dr. Supervisor Name (Thesis Supervisor);
        \item Prof. Dr. External Expert Name (External Expert).
    \end{itemize}
}

\supervisor{Prof. Dr. Supervisor Name}                % name of supervisor
\dean{Prof. Dr. Dean Name}                            % name of faculty dean

\acknowledgments{I would like to thank my supervisor, colleagues, and family for their continuous support throughout this research project.}


%----------------------------------------------------------------------------------------
%	BEGIN DOCUMENT AND CREATE TITLEPAGE
%----------------------------------------------------------------------------------------
\maketitle


%----------------------------------------------------------------------------------------
%	PREAMBLE
%----------------------------------------------------------------------------------------
\begin{abstractstyle}{\hsummary}
    This thesis investigates \textit{[research topic]} in the context of \textit{[application domain]}. The study addresses the core problem of \textit{[problem statement]}, which remains insufficiently resolved due to \textit{[gap 1]} and \textit{[gap 2]}. This limitation motivates the need for a more reliable and practically applicable solution.

    To address this problem, the research adopts a design-develop-evaluate approach. The methodology combines \textit{[data sources or context]}, \textit{[methods or architecture]}, and \textit{[evaluation strategy]} to ensure methodological rigor and real-world relevance. The proposed approach is designed to align with the stated research questions and objectives.

    The study delivers \textit{[artifact or system outcome]} and evaluates its performance against \textit{[baseline or benchmark criteria]}. The findings are expected to contribute to \textit{[research-domain contribution]} and \textit{[problem-domain contribution]}, while also identifying key limitations and future directions for extended investigation.

    \textbf{ACM Subject Descriptors:} \textit{[descriptor 1]}; \textit{[descriptor 2]}; \textit{[descriptor 3]}.

    \textbf{Keywords:} \textit{[keyword 1]}; \textit{[keyword 2]}; \textit{[keyword 3]}; \textit{[keyword 4]}; \textit{[keyword 5]}.
\end{abstractstyle}

\thesisprintdeclaration
\thesisprintacknowledgments


% table of contents ----------------------------------------------------------------------

\tableofcontents
\listoffigures
\listoftables
% print list of acronyms and glossary
\printnoidxglossaries



%----------------------------------------------------------------------------------------
%	MAIN CONTENT
%----------------------------------------------------------------------------------------
\mainmatter

\chapter{Introduction}
\section{Chapter overview}
State the flow of this chapter using the TTT approach: context, problem definition, motivation, gap, contributions, challenges, research questions, aim, and objectives.

\section{Problem background}
\subsection{Context of the research domain}

Analysing company shares has become an increasingly complex and expertise-driven task in modern financial markets. Effective evaluation requires investors to review numerous detailed documents such as annual reports, quarterly statements, regulatory announcements, and news articles, many of which contain highly technical information. Interpreting these materials demands a solid understanding of accounting and finance, skills that many individual investors lack. Even for those equipped with the necessary knowledge, analysing many interrelated parameters mentally remains a challenging and time-consuming process. 

Another major challenge arises from the contextual nature of financial markets. Analytical models developed for one country’s market often fail to produce accurate results when applied to another. Market behaviour is influenced by a range of country-specific factors, including economic conditions, availability of resources, governance structures, and political stability. Therefore, analytical models must be adapted to reflect the unique characteristics of each national market. 

Given these complexities, most high-net-worth individuals and brokerage firms employ professional analysts to conduct market evaluations manually. In contrast, retail investors often lack the time, expertise, and resources to perform such detailed analysis. As a result, they become more susceptible to misinformation, market manipulation, and financial losses. 

Retail investors frequently rely on research and recommendations provided by brokerage firms. However, over time, it has become evident that such research cannot always be considered objective or reliable. Brokerage firms often face inherent conflicts of interest, as their revenue tends to increase with higher trading volumes. Consequently, investors may be encouraged to trade more frequently, even when such activity is not in their best financial interest. 

This imbalance in knowledge and analytical capacity contributes to market inefficiencies. Due to the limited financial literacy within the investor community, certain companies may trade at unrealistically high valuations, while others remain undervalued conditions that are detrimental to both market health and broader economic stability. 

Furthermore, brokerage firms often issue target prices for companies, but forecasting such figures is inherently uncertain and difficult. Effective share market analysis should instead prioritize a deep understanding of the underlying business fundamentals rather than relying solely on externally set price targets. Before projecting future performance, it is essential to examine a company’s historical financial data and assess its current operational status. 

Although there has been substantial attention on technical analysis, the development and accessibility of robust fundamental analysis models remain limited. 

\subsection{Subdomain and key technological context}
Narrow the focus to the specific technical area that your project addresses.

\subsection{Practical and real-world context}
Explain application context, adoption constraints, and current industry or regulatory relevance.

\section{Problem definition}

\subsection{Problem statement}
Fundamental analysis is widely used in practice but remains relatively under-researched, primarily due to its reliance on large volumes of unstructured data from diverse sources. This makes it difficult for human analysts to systematically process and interpret information, highlighting the need for an Agentic AI-based approach to enable more efficient and consistent analysis. 

\section{Research motivation}

The analysis of equity investments in the Sri Lankan capital market remains a complex and often inaccessible task for many investors, particularly those with limited financial domain knowledge. At present, investors primarily rely on general-purpose large language model (LLM) tools and research reports provided by broker firms. While these resources offer some level of guidance, they are not specifically tailored to the nuances of stock market analysis, resulting in inefficiencies and a steep learning curve for users. 

Consequently, a significant number of investors are prone to misinterpretation of information and suboptimal decision-making, leading to potential financial losses. This challenge is especially pronounced among retail investors, who may lack the expertise required to critically evaluate financial data and insights. Moreover, these limitations can negatively impact overall confidence and trust in the Sri Lankan capital market. 

Therefore, there is a clear need for a more specialized, accessible, and intelligent solution that can support investors in making informed decisions, while also enhancing transparency and trustworthiness within the market ecosystem. 

\section{Research gap}
\textbf{Gap 1:} Add citation-backed explanation.

\textbf{Gap 2:} Add citation-backed explanation.

\section{Contribution to the body of knowledge}
\subsection{Contribution to the research domain}
Explain methodological or theoretical contribution.

\subsection{Contribution to the problem domain}
Explain practical or applied contribution.

\section{Research challenges}
\textbf{Challenge 1:} Describe challenge and justification.

\textbf{Challenge 2:} Describe challenge and justification.

\section{Research questions}
\resq{1}{What key parameters can be extracted using AI to improve the efficiency and effectiveness of fundamental analysis?}
\resq{2}{What is the influential impact (in detail and color context) of a multi-scale feature
extraction mechanism on enhancing low-light images in non-uniform lighting conditions?}
\resq{3}{In what ways can the GAN-based networks be directed to produce robust and generalized
enhancements across a diverse non-uniformity range of low-light scenes?}

\section{Research aim}
State the aim in one sentence using design-develop-evaluate language, then elaborate briefly.

\section{Research objectives}
\begin{itemize}
    \item RO1: Insert SMART objective 1.
    \item RO2: Insert SMART objective 2.
    \item RO3: Insert SMART objective 3.
    \item RO4: Insert SMART objective 4.
\end{itemize}

\section{Chapter summary}
Summarize key points and transition to Chapter 2.

\chapter{Literature review}
\section{Chapter overview}
Describe how this chapter reviews concepts, technologies, existing work, and evaluation practices.

\section{Concept map}
Insert a conceptual map figure and provide a brief explanation of concept relationships.

\begin{figure}[H]
    \centering
    \fbox{\rule{0pt}{2in}\rule{0.9\linewidth}{0pt}}
    \caption{Concept map of the research domain (placeholder)}
    \label{fig:concept-map}
\end{figure}

\section{Problem domain review}
\subsection{Introduction to the domain}
Define the overarching domain and its key terms.

\subsection{Key challenges and sub-areas}
Review major thematic challenges relevant to your problem.

\subsection{Existing frameworks, standards, and domain approaches}
Summarize applicable standards, regulations, and best-practice frameworks.

\subsection{Proposed architecture positioning}
Position your high-level architecture against the reviewed domain context.

\begin{figure}[H]
    \centering
    \fbox{\rule{0pt}{2in}\rule{0.9\linewidth}{0pt}}
    \caption{Proposed conceptual architecture (placeholder)}
    \label{fig:proposed-architecture}
\end{figure}

\section{Technological review}
\subsection{Data acquisition and preprocessing}
Critically review data handling techniques used in comparable work.

\subsection{Model or system architectures}
Compare relevant architecture choices and trade-offs.

\subsection{Tuning and optimization approaches}
Review methods used to optimize performance and robustness.

\subsection{Evaluation metrics used in prior work}
Evaluate suitability of common metrics for your project context.

\section{Existing work}
\subsection{Thematic grouping}
Group studies by themes and compare their findings.

\subsection{Cross-theme linking}
Explain overlaps and differences across themes and methods.

\section{Evaluation and benchmarking}
Critically analyze datasets, protocols, metrics, and evaluation limitations.

\section{Chapter summary}
Synthesize key gaps and connect them to Chapter 3.

\chapter{Methodology}
% \section{Chapter overview}
% Outline the methodological structure and rationale.

\section{Research methodology}

% The following table provides an overview of the selected scientific methodology based on the
% conventional Research ‘Onion’ Model, presented by (Saunders et al., 2019).

% Table

This study adopts pragmatism as its research philosophy, as it is primarily focused on addressing a real-world problem, enhancing fundamental analysis through an Agentic AI-based approach, rather than adhering to a single philosophical paradigm. Pragmatism enables the integration of both quantitative methods and qualitative insights, allowing the research to effectively handle diverse data types. 

\section{Development methodology}

This study adopts an iterative development methodology, where the proposed Agentic AI-based system is developed in incremental stages. Each iteration involves designing and implementing a prototype, followed by systematic evaluation of its performance in supporting fundamental analysis. The insights gained from each evaluation are used to refine and enhance the subsequent version of the system. This continuous cycle of development and improvement ensures that the solution evolves progressively, leading to a more effective and reliable framework for fundamental analysis. 

\section{Project management methodology}

Agile principles are adopted for project management to support iterative development and continuous improvement, enabling the system to evolve through successive refinements until it meets the desired project objectives.

\subsection{Project Scope}

This study focuses on fundamentally analysing publicly listed companies in the Colombo Stock Exchange (CSE). However, diversified holding companies are excluded due to the complexity arising from their multi-sector operations. The analysis is primarily based on annual reports, and quarterly reports are not considered. As a result, the findings may reflect a maximum time lag of up to one year. 

\subsection{Schedule}

\begin{table}[H]
    \centering
    \begin{tabular}{|l|l|}
        \hline
        Milestone & Target date \\
        \hline
        Project Proposal Presentation  & May 7th, 2026  \\
        Project Proposal and Interim Demonstration  & May 14th, 2026  \\
        Implementation complete & June 30th, 2026  \\
        Submission  & July 30th, 2026 \\
        \hline
    \end{tabular}
    \caption{Project schedule summary}
    \label{tab:project-schedule}
\end{table}

% \subsection{Deliverables and dates}
% List key deliverables and deadlines.

% \subsection{Resources}
% \subsubsection{Hardware resources}
% List compute and infrastructure resources.

% \subsubsection{Software resources}
% List tools, frameworks, and libraries.

% \subsubsection{New technical skills}
% List technical skills acquired during the project.

% \subsection{Data requirements}
% Specify datasets, licensing, and preprocessing requirements.

% \subsection{Risks and mitigation}
% Document major project risks with likelihood, impact, and mitigation strategy.

% \begin{table}[H]
%     \centering
%     \begin{tabular}{|l|l|l|l|}
%         \hline
%         Risk & Likelihood & Impact & Mitigation \\
%         \hline
%         Data quality risk & Medium & High & Validation and cleaning pipeline \\
%         Timeline overrun & Medium & Medium & Weekly milestone tracking \\
%         \hline
%     \end{tabular}
%     \caption{Risk and mitigation table (placeholder)}
%     \label{tab:risk-mitigation}
% \end{table}

% \section{Chapter summary}
% Summarize methodological choices and transition to LESP and implementation.

\chapter{LEGAL, ETHICAL, SOCIAL AND PROFESSIONAL ISSUES}
\section{Chapter overview}
Introduce why LESP considerations are critical for the research context.

\section{BCS Code of Conduct}
Identify relevant BCS principles and map them to your project decisions.

\section{Social issues}
Discuss bias, accessibility, and wider societal impacts.

\section{Legal issues}
Discuss data protection, consent, and intellectual property compliance.

\section{Ethical issues}
Discuss anonymization, fairness, transparency, and accountability.

\section{Professional issues}
Discuss responsible reporting, professional integrity, and disclosure practices.

\section{Chapter summary}
Summarize key LESP considerations and mitigations.

\chapter{IMPLEMENTATION AND EXPERIMENTAL SETUP}
\section{Chapter overview}
Summarize the implementation flow and what this chapter covers.

\section{Technology selection}
\subsection{Programming language}
State language choice and technical justification.

\subsection{Libraries}
List core libraries and justify their roles.

\subsection{Frameworks}
List frameworks and architectural role.

\subsection{IDEs and tools}
List development, collaboration, and testing tools.

\begin{table}[H]
    \centering
    \begin{tabular}{|l|l|l|l|}
        \hline
        Component & Name & Purpose & Justification \\
        \hline
        Language & Insert language & Core development & Insert rationale \\
        Framework & Insert framework & Service layer & Insert rationale \\
        Library & Insert library & Core functionality & Insert rationale \\
        \hline
    \end{tabular}
    \caption{Technology stack summary (placeholder)}
    \label{tab:technology-stack}
\end{table}

\section{Data selection}
Describe selected data sources, selection rationale, licensing, and preprocessing summary.

\section{Implementation of core functionalities}
\subsection{Module 1}
Describe module purpose and implementation notes.

\subsection{Module 2}
Describe module purpose and implementation notes.

\subsection{Module 3}
Describe module purpose and implementation notes.

\section{User interface (if applicable)}
Describe implemented UI elements and usability considerations.

\begin{figure}[H]
    \centering
    \fbox{\rule{0pt}{2in}\rule{0.9\linewidth}{0pt}}
    \caption{Sample user interface screenshot (placeholder)}
    \label{fig:ui-screenshot}
\end{figure}

\section{Chapter summary}
Summarize implemented components and transition to evaluation.

\chapter{Results and evaluation}
\section{Chapter overview}
Describe the evaluation types and how they validate your research outcomes.

\section{Quantitative results}
Present numerical metrics with clear explanation of trends and anomalies.

\begin{table}[H]
    \centering
    \begin{tabular}{|l|l|l|l|}
        \hline
        Model & Accuracy & F1 score & Dataset \\
        \hline
        Baseline & Insert value & Insert value & Insert dataset \\
        Proposed & Insert value & Insert value & Insert dataset \\
        \hline
    \end{tabular}
    \caption{Quantitative performance metrics (placeholder)}
    \label{tab:quantitative-results}
\end{table}

\section{Qualitative analysis}
Interpret visual outputs, case studies, and domain-expert observations.

\section{Ablation or sensitivity studies}
Analyze the impact of removing or changing major components.

\begin{table}[H]
    \centering
    \begin{tabular}{|l|l|l|}
        \hline
        Variant & F1 score & Remarks \\
        \hline
        Full model & Insert value & Reference setup \\
        Variant A & Insert value & Effect of change \\
        \hline
    \end{tabular}
    \caption{Ablation study results (placeholder)}
    \label{tab:ablation-results}
\end{table}

\section{Baseline or SOTA comparison}
Compare proposed approach with baselines and state-of-the-art methods using fair protocols.

\section{Discussion}
Interpret findings against research questions, objectives, and practical implications.

\section{Chapter summary}
Summarize key findings and transition to conclusion and future work.

\chapter{CONCLUSION AND FUTURE WORK}
\section{Chapter overview}
Describe how this chapter consolidates findings and outlines future directions.

\section{Summary of findings}
Map each objective and research question to achieved outcomes.

\begin{table}[H]
    \centering
    \begin{tabular}{|l|l|l|}
        \hline
        Objective & Achievement & Evidence chapter \\
        \hline
        RO1 & Insert achievement & Chapter 2 \\
        RO2 & Insert achievement & Chapters 3-4 \\
        RO3 & Insert achievement & Chapter 5 \\
        RO4 & Insert achievement & Chapter 6 \\
        \hline
    \end{tabular}
    \caption{Objectives versus achievements (placeholder)}
    \label{tab:objectives-achievements}
\end{table}

\section{Contributions}
\subsection{Contribution to the research domain}
Re-state theoretical or methodological advancements.

\subsection{Contribution to the problem domain}
Re-state practical and applied impact.

\section{Limitations}
Document data, technical, evaluation, and contextual limitations.

\section{Future work}
List specific, realistic follow-up research directions derived from limitations.

\section{Contribution to the body of knowledge}
Explicitly articulate how completed work advances both domains.

\section{Concluding remarks}
Provide a concise final reflection and forward-looking statement.


\bibliographystyle{apalike}
\footnotesize\bibliography{references}



%----------------------------------------------------------------------------------------
%	APPENDIX
%----------------------------------------------------------------------------------------
\appendix

\chapter{Appendix overview}
Use appendices for supplementary material only. Keep all critical arguments, methodology, and results in the main chapters.

\section{Writing sequence checklist}
\begin{enumerate}
    \item Finalize Chapter 1 (problem, gap, research questions, aim, objectives).
    \item Complete Chapter 2 based on Chapter 1 scope and research questions.
    \item Finalize Chapter 3 and Chapter 4 before implementation details.
    \item Develop Chapter 5 and Chapter 6 iteratively as implementation and experiments progress.
    \item Write Chapter 7 after final evaluation outcomes are stable.
\end{enumerate}

\section{Final quality gate checklist}
\begin{itemize}
    \item Verify numbering flow: cover page unnumbered, front matter roman, main chapters arabic, appendix separate.
    \item Ensure every figure and table has a caption and explanatory reference in text.
    \item Ensure each in-text citation is present in the references list and vice versa.
    \item Verify all research questions and objectives are explicitly answered by the end of Chapter 7.
    \item Confirm critical results remain in main chapters and appendices hold supplementary material only.
\end{itemize}

\section{Detailed concept map}
Insert the full-size concept map used in Chapter 2.

\section{Detailed project schedule}
Insert the full Gantt chart or equivalent timeline artifact.

\section{Additional tables and figures}
Insert large supporting tables, extended results, and non-critical supplementary visuals.

\section{Code snippets and technical artifacts}
Insert selected supplementary code snippets and implementation artifacts if required.


\glsaddallunused                                % add all unused items to glossary

\end{document}
